% ============================================================================
% Itérations du projet HackathonQuantique : couleurs et marqueurs communs.
% Inclus par tous les cours et par la fiche projet, pour qu'une itération ait
% la même couleur partout. Pour ajouter une itération : une couleur, un nom,
% une date, une ligne dans \legendeiterations.
% ============================================================================
\definecolor{itV0}{HTML}{8D6E63}   % brun   : V0, modèle de la note d'intention
\definecolor{itV1}{HTML}{C2185B}   % framboise : V1, chaîne complète conforme à la note
\definecolor{itV2}{HTML}{00838F}   % cyan foncé : V2, problème reformulé
\definecolor{itV3}{HTML}{E09A00}   % ambre  : V3, concurrent équitable et problème plus grand
\definecolor{itV4}{HTML}{5E35B1}   % indigo : V4, réservé

\makeatletter
\newcommand{\itername}[1]{\@nameuse{itname@#1}}
\newcommand{\iterdate}[1]{\@nameuse{itdate@#1}}
\newcommand{\defiter}[3]{\@namedef{itname@#1}{#2}\@namedef{itdate@#1}{#3}}
\makeatother
\defiter{V0}{Modèle de la note d'intention}{24--30 oct. 2026}
\defiter{V1}{Chaîne complète conforme à la note}{4--5 nov. 2026}
\defiter{V2}{Problème reformulé}{5--6 nov. 2026}
\defiter{V3}{Concurrent équitable, problème plus grand}{à partir du 6 nov. 2026}

% Pastille : \iterbadge{V1}
\newcommand{\iterbadge}[1]{\tikz[baseline=(b.base)]\node[fill=it#1,text=white,rounded corners=2pt,
  inner xsep=3pt,inner ysep=1.5pt,font=\sffamily\bfseries\scriptsize] (b) {#1};}

% Bandeau de début de chapitre : \iterbanner{V0}
\newcommand{\iterbanner}[1]{\par\noindent\begin{tikzpicture}
  \node[fill=it#1!12,draw=it#1,line width=0.6pt,rounded corners=2pt,text width=\dimexpr\linewidth-12pt\relax,
    inner sep=4pt,anchor=west,font=\small] at (0,0)
    {\iterbadge{#1}\ \textbf{\color{it#1}Itération #1} --- \itername{#1} \hfill{\color{it#1}\iterdate{#1}}};
  \end{tikzpicture}\par\medskip}

% Encadré d'évolution : \begin{evolution}{V1}{titre} ... \end{evolution}
\newtcolorbox{evolution}[2]{enhanced,breakable,boxrule=0pt,leftrule=3pt,arc=0pt,
  left=6pt,right=6pt,top=4pt,bottom=4pt,colback=it#1!7,colframe=it#1,
  fonttitle=\bfseries\small,coltitle=it#1,attach title to upper={\ },
  title={\iterbadge{#1}\ Évolution #1 : #2.}}

% Légende des itérations, pour les avant-propos
\newcommand{\legendeiterations}{%
\begin{center}\small
\begin{tabular}{lll}
\iterbadge{V0} & \itername{V0} & \iterdate{V0}\\
\iterbadge{V1} & \itername{V1} & \iterdate{V1}\\
\iterbadge{V2} & \itername{V2} & \iterdate{V2}\\
\iterbadge{V3} & \itername{V3} & \iterdate{V3}\\
\end{tabular}
\end{center}}

% Encadré « pistes d'amélioration » : défini ici s'il manque (cours V0 sans cet encadré)
\providecolor{cviolet}{HTML}{7A4FB3}
\ifcsname pistes\endcsname\else
\newtcolorbox{pistes}[1][]{enhanced,breakable,boxrule=0pt,leftrule=2.5pt,arc=0pt,
  left=6pt,right=6pt,top=4pt,bottom=4pt,fonttitle=\bfseries\small,coltitle=cink,attach title to upper={\ },
  colback=cviolet!6,colframe=cviolet,title={Pistes d'amélioration#1.}}
\fi
